\par\Needspace{6\baselineskip}\medskip
\begin{center}\hypertarget{native-246-article175}{}\pdfbookmark[2]{175.}{article-175}175.\end{center}
To this end we will pass from our particular problem to a very general inquiry, most fruitful in every. application of calculation to natural philosophy. Let $V$, $V'$, $V''$ etc. be functions of the unknowns $p$, $q$, $r$, $s$ etc., let $\mu$ be the number of those functions and $\nu$ the number of unknowns, and suppose that direct observations have yielded the following values of the functions: $V=M$, $V'=M'$, $V''=M''$ etc. Generally speaking, finding the values of the unknowns will thus be an indeterminate, determinate, or overdetermined problem, according as $\mu<\nu$, $\mu=\nu$, or $\mu>\nu$\footnote{If in the third case the functions $V$, $V'$, $V''$ etc. were such that $\mu+1-\nu$ or more of them could be regarded as functions of the others, the problem would still be overdetermined with respect to these functions, but indeterminate with respect to the quantities $p$, $q$, $r$, $s$ etc.: their values could not be determined even if absolutely exact values of the functions $V$, $V'$, $V''$ etc. were given: but we shall exclude this case from our inquiry.}.
We shall discuss only the last case, in which clearly an exact representation of all the observations would be possible only if they were all absolutely free of errors. Since this does not happen in nature, every system of values of the unknowns $p$, $q$, $r$, $s$ etc. must be regarded as possible if it yields values of the \SourceLeaf{247}functions $M-V$, $M'-V'$, $M''-V''$ etc. no greater than the limits of the errors that could have been made in these observations: this must by no means be understood to imply that these possible systems all have the same degree of probability.

Let us first suppose that the circumstances were such in all the observations that there is no reason to suspect that one is less exact than another, or, in other words, that errors of the same size are to be regarded as equally probable in each. The probability to be assigned to any error $\Delta$ will accordingly be expressed by a function of $\Delta$ itself, which we shall denote by $\varphi\Delta$. Even if we cannot specify this function precisely, we can at least say that its value must be greatest for $\Delta=0$, that it will generally be the same for values of $\Delta$ equal in magnitude and opposite in sign, and finally that it vanishes if $\Delta$ is taken as the greatest error or a greater value. Strictly speaking, therefore, $\varphi\Delta$ must be classed among the discontinuous functions, and if for practical purposes we allow ourselves to substitute an analytic function for it, that function must be such that on either side of $\Delta=0$ it converges, as it were asymptotically, to $0$, so that beyond that limit it may be regarded as actually vanishing. Further, the probability that the error lies between the limits $\Delta$ and $\Delta+\mathrm{d}\Delta$, which are separated by the infinitely small difference $\mathrm{d}\Delta$, is to be expressed by $\varphi\Delta\, .\,\mathrm{d}\Delta$; hence in general the probability that the error lies between $D$ and $D'$ will be given by the integral $\int\varphi\Delta\, .\,\mathrm{d}\Delta$, taken from $\Delta=D$ to $\Delta=D'$. This integral, taken from the extreme negative value of $\Delta$ to its greatest positive value, or more generally from $\Delta=-\infty$ to $\Delta=+\infty$, must necessarily be $=1$.

Suppose, then, that a certain determinate system of values of the quantities $p$, $q$, $r$, $s$ etc. holds. The probability that the observation will yield the value of $V$ as $M$ will be expressed by $\varphi(M-V)$, once we have substituted in $V$ the corresponding values for $p$, $q$, $r$, $s$ etc.; similarly $\varphi(M'-V')$, $\varphi(M''-V'')$ etc. will express the probabilities that the observations will yield for the functions $V'$, $V''$ etc. the values $M'$, $M''$ etc. Accordingly, since all the observations may be regarded as independent events, the product
\[
\varphi(M-V)\, .\,\varphi(M'-V')\, .\,\varphi(M''-V'')\, .\,\text{etc.}=\Omega
\]
will express the expectation or probability that all these values will result from the observations together.
